\subsection{Implementation Specifics of Scenarios}\label{subsec:implementation-specifics-of-scenarios}
This section describes scenario-specific implementation details built on top of the framework.

\subsubsection{Corporate}
The goal is to implement website for a software agency, that advertise their services, projects, and blog articles
using the CMS features.

The Home directory provided by the framework does not include JavaScript file,
but the project homepage requires extra JavaScript functionality and thus home.js is created
and the functionality is implemented.
The template is edited to import the script file using SideLoader's Javascript::import() method and the homepage layout
is designed and implemented.

The framework currently does not support e-commerce features, which limits the usefulness in corporate environment.
However, simple corporate requirements can be met with initial installment of the framework.

\subsubsection{Content-Driven}
The goal is to implement website for the administrators to share a curated repository of recipes and cooking tips
using the articles and role-based privilege systems.

The homepage needs to be updated to render last edited articles and list category pages.
The framework includes all the core functionalities needed to implement the behavior.
The following functions are implemented for the Home component to be accessible in the template.

\begin{verbatim}
    public function getCategoryPage(): ?Page {
        // Get from admin settings category parent page id
        $pageId = Setting::fromName(
            'savora:homepage-categories-page-id',
            true,
            1,
            [PROPERTY_EDITABLE => true]
        )->toInt();

        return Page::fromId($pageId);
    }

    public function getFeaturedPages(): array {
        // Get from admin settings how many articles to show
        $limit = Setting::fromName(
            'savora:homepage-featured-pages-count',
            true,
            3,
            [PROPERTY_EDITABLE => true]
        )->toInt();

        return Page::lastUpdated(limit: $limit);
    }
\end{verbatim}

With all needed backend features completed the template can be edited with custom homepage design.

In the administration dashboard, the user group Editor is created and privileges for content management are granted.
The group contains only one test user with the tag \("fizz"\),
which is a system-wide unique identifier used for authentication, and password \("fizzbuzz"\).

The framework currently supports only closed systems, in which the administrators create users who can manage the content.
This decision introduces a trade-off: it simplifies the implementation of content-driven systems, but reduces flexibility.

\subsubsection{Personal}
The goal is to create a Chrome extension which shows static image as background and displays widgets editable with built-in editor.
To achieve this goal, the core of the extension is interactive webpage.
This enables live updates for the extension using HTML iframe element to load it, and up-to-date demonstration.
This architectural design introduces a trade-off: the extension requires a stable internet connection because the interface is remotely hosted.
However, after loading the interface, all content is generated from configuration stored directly in the user's browser.

The extension content page is part of the Action and View framework's systems, which enable direct bind to URL path.

\begin{verbatim}
    // index.php
    // Frame is the extension content page
    // Bound to protocol://domain/extension
    $router->use('/extension', new Frame());
\end{verbatim}

The page can be added to pages in administration as external page using fully qualified URL\@.

The freeform widget editor needs to be implemented by hand.
Support for custom editor is currently limited, still the framework provides general tools.
This includes asset loading, general JavaScript utilities, forms, inspector controls from webpage content widget editor.

The scenario implements serving images from selected server directories as extra feature.
To complete the feature web settings in administration settings are used and custom endpoint with following workflow:

\begin{enumerate}
    \item Get background defined in web settings
    \item Compute file difference in defined directories from last request
    \item Update internal files listing if any of the directories have been updated
    \item Select random image from internal files listing
    \item Send URL to access the selected image
\end{enumerate}

On extension startup, the endpoint is called and if URL to image is returned it is used as background image.
To manage directories the administrator must navigate to System \(>\) Settings and edit
the \("\)startuh:background\_directory\_os\_hash\("\) setting.
Multiple directories are supported by using the \(;\) directory delimiter.

The personal scenario demonstrates that the framework can be extended beyond traditional CMS-based applications.
Although interactive elements require manual implementation.
